\documentclass{article}
\usepackage[T1]{fontenc}
\usepackage{lmodern}
\usepackage{graphicx}
\usepackage{amsmath,amssymb}
\usepackage[a4paper,margin=2cm]{geometry}
\usepackage[absolute,overlay]{textpos}
\setlength{\TPHorizModule}{1mm}
\setlength{\TPVertModule}{1mm}
\usepackage[indonesian]{babel}
\usepackage{hyperref}
\setlength{\emergencystretch}{2em}
% unit: Halaman 9

\begin{document}

% === Halaman 9 (llm) ===

\section*{Medan (\textit{Field})}

\begin{itemize}
    \item Medan (\textit{field}) adalah system aljabar yang terdiri dari
    \begin{itemize}
        \item himpunan elemen (disimbolkan dengan F)
        \item operasi biner, yaitu penjumlahan (+) dan perkalian ($\cdot$).
    \end{itemize}

    \item Sebuah struktur aljabar $\langle F, +, \cdot \rangle$ disebut medan jika dan hanya jika:
    \begin{enumerate}
        \item $\langle F, + \rangle$ adalah grup abelian
        \item $\langle F - \{0\}, \cdot \rangle$ adalah grup abelian
        \item Operasi $\cdot$ menyebar terhadap operasi $+$ (sifat distributif)
    \end{enumerate}
    Distributif:
    \[
    \begin{aligned}
    x \cdot (y + z) &= (x \cdot y) + (x \cdot z) \\
    (x + y) \cdot z &= (x \cdot z) + (y \cdot z)
    \end{aligned}
    \]

    \item Jadi, sebuah medan memenuhi aksioma: \textit{closure}, komutatif, asosiatif, dan distributif
\end{itemize}

\end{document}
